\section{Positive wall transport and change of chart}
\label{wall:section}

The reordering theorem gives positive local changes of coordinates.  We
now assemble them into a consistent system, define its incoming
coordinates, and prove the two properties needed later: positivity of
sections and compatibility with a change of simple roots.  All wall
arrangements in the argument are finite after a specified truncation.

This is a scattering-diagram construction in the framework originating
with Kontsevich--Soibelman and Gross--Siebert \cite{KS06,GS11}.  Its
description by a total transport element is the wall-crossing
construction of \cite[Theorem~2.1.6 and Proposition~2.1.12]{KS14},
expressed here in completed quantum-torus coordinates; the locators
refer to the version recorded in the bibliography.  Quantum wall and
theta positivity are developed in \cite{DM21}.  We retain direct
proofs to specify exactly the support bounds and coefficient algorithms
used by the permutation map.

\subsection{The completion and the incoming prescription}

Let $M$ be a lattice with an integral alternating form $\Omega$, and put
$K=\Q((v^{-1}))$.  We use the centered quantum torus convention
\begin{equation}
 X_aX_b=v^{\Omega(a,b)}X_{a+b},
 \qquad V(a)=\Omega(\,\cdot\,,a)\in M^\vee_{\R}.
 \label{wall:torus}
\end{equation}
We may assume that $\Omega$ is nondegenerate: replace $M$ by
$M\oplus M^\vee$ and extend the form by
$\Omega((a,f),(b,g))=\Omega(a,b)+f(b)-g(a)$.
The original monomials form a subalgebra, so this enlargement changes
none of their computations.

A \emph{chart} is a finite set $C\subset M$ contained in a lattice
basis.  Write
\[
 C^+=\sum_{c\in C}\N c,\qquad
 |\textstyle\sum_c u_cc|_C=\sum_cu_c,\qquad
 a\le_C b\ \Longleftrightarrow\ b-a\in C^+.
\]
Every interval $\{a:m\le_C a\le_C n\}$ is finite.  The completed
positive group is
\[
 G_C=1+\prod_{r\in C^+\setminus\{0\}}KX_r.
\]
Multiplication, inversion, logarithms, and exponentials are defined
by root degree.  A coefficient of any such operation is a finite sum
over decompositions of its exponent.  A section value will belong to
the corresponding completed torus module supported in a finite union
of translates of $C^+$.  These support conditions are maintained
throughout.

For a covector $h$, let $G_{C,+}(h)$, $G_{C,0}(h)$, and $G_{C,-}(h)$
denote the subgroups supported on roots with, respectively, positive,
zero, and negative $h$-value.  Every $g\in G_C$ has a unique
factorization
\begin{equation}
 g=g_+(h)g_0(h)g_-(h).
 \label{wall:three-factors}
\end{equation}
Indeed, suppose all coefficients below degree $d$ have been found.
For a root $r$ of degree $d$, its coefficient in the product is the
as-yet unknown coefficient in exactly one of the three factors, plus
known products of smaller roots.  Subtraction determines that
coefficient.  This also proves the following refinement rule: near
$h$, the factors are obtained by factoring $g_0(h)$ according to the
perturbation of $h$, while keeping the already strict signs fixed.
Both assertions may equally be made modulo roots outside a finite
addition-order ideal.

The factors define walls.  At a generic point of $r^\perp$, the
middle factor is supported on the positive ray through $r$; it may
be the identity on some portions of that hyperplane.  Crossing from
positive to negative on that ray applies
$\Ad(g_0):F\mapsto g_0Fg_0^{-1}$, and crossing in reverse applies its
inverse.  Products of crossing elements have the earlier element on
the right.  The transformation from an all-positive chamber to $h$
is $\Ad(g_-(h))$. Indeed, for a positive-to-negative crossing from
$I$ to $II$, the refinement rule gives
\[
 g_-(II)=g_0g_-(I),\qquad
 \Ad(g_-(II))=\Ad(g_0)\circ\Ad(g_-(I)).
\]
Thus each crossing agrees with the transformation determined by its
endpoint, proving path independence and the stated order of composition.
At degree at most $d$, there are finitely many roots, hyperplanes,
and cells, so generic polygonal paths and their crossings have their
usual literal meaning at that order.  We use no assertion of local
finiteness for the untruncated arrangement.

Conversely, a consistent assignment of such wall elements, constant
on the cells of finite arrangements at each order, comes from
\eqref{wall:three-factors}.  Take $g$ to be the ordered product from
an all-positive to an all-negative chamber.  Near a given generic
crossing point, a perturbed straight path to that point crosses only
roots negative there; a path completing the journey crosses only
roots positive there.  The resulting factorization of $g$ has the
prescribed crossing element as its middle factor, by uniqueness.

Each simple root carries the specific unit
\begin{equation}
 D_p=\prod_{j\ge0}(1+v^{-1-2j}X_p),\qquad p\in C.
 \label{wall:initial-unit}
\end{equation}
Define the wall system by requiring
\begin{equation}
 [X_r]\log g_0(V(r))=
 \begin{cases}
 [X_r]\log D_p,&r=jp,\quad p\in C,\ j\ge1,\\
 0,&\text{otherwise}.
 \end{cases}
 \label{wall:incoming-prescription}
\end{equation}
These equations determine $g$ uniquely.  To see this directly, write
$\log g=\sum_r\ell_rX_r$ and solve by $|r|_C$.  In factoring at
$V(r)$, the linear term $\ell_rX_r$ belongs to the middle factor,
since $V(r)(r)=0$.  All other contributions to its coefficient use
strictly smaller roots.  Thus the equation for $r$ has coefficient
one on $\ell_r$ and determines it from previous equations.  In
particular, every bounded-order part of the system is computable.

Here and below, ``incoming'' refers to the end of a wall in the
direction $V(r)$.  We verify that this geometric meaning agrees with
\eqref{wall:incoming-prescription}, even when $V(r)$ lies on several
hyperplanes.  On the algebra supported in $\ker V(r)$, the functional
$\tau_r(B)=[X_r]B$ is a trace: a term with exponents $a+b=r$ has
$\Omega(a,b)=\Omega(a,r)=0$.  Consequently
\begin{equation}
 \tau_r\log(B_1B_2)=\tau_r\log B_1+\tau_r\log B_2
 \label{wall:trace}
\end{equation}
for elements with constant term one in that algebra.  For completeness,
replace each $X_a$ by $z^{|a|_C}X_a$ and differentiate.  Cyclic
permutation in the trace gives
$\tau_r((\log B)')=\tau_r(B'B^{-1})$; the product rule proves equality
of the derivatives in \eqref{wall:trace}, and the constant terms are
zero.  This is a formal calculation at each root degree.
Factor $g_0(V(r))$ using a small generic displacement within
$r^\perp$.  Of those factors, only the middle one can contain $X_r$.
Equation~\eqref{wall:trace} therefore identifies its logarithmic
$r$-coefficient with the prescribed one.  Applying the argument to
each positive multiple of the primitive ray shows that the incoming
wall element is exactly $D_p$ on a simple ray and the identity on
every other ray.

\subsection{Positive wall factors and incoming coordinates}

Recall the two elementary unit types from Theorem~\ref{alg:reordering}:
\[
 U_{a,k,f}=\prod_{j\ge0}f(v^{k-2j}X_a),\qquad
 f(T)=1+T\quad\hbox{or}\quad f(T)=(1-T)^{-1}.
\]
An allowed list has its parameters $k$ bounded above at bounded root
degree, and has only finitely many occurrences above any fixed lower
threshold on $k$.  Occurrences, including repeated units, are counted
separately.

\begin{lemma}[The unit normalization and positive crossing]
\label{wall:unit}
For $j\ge0$,
\begin{equation}
 [X_{jp}]D_p=
 \frac{v^{-j^2}}{\prod_{i=1}^j(1-v^{-2i})}.
 \label{wall:unit-coefficient}
\end{equation}
Let $t=\Omega(a,m)$.  Then
\begin{equation}
 \Ad(U_{a,k,f})X_m
 =X_m\frac{U_{a,k,f}(v^{2t}X_a)}{U_{a,k,f}(X_a)}.
 \label{wall:unit-ratio}
\end{equation}
For $t>0$ the ratio is $\prod_{\ell=1}^t f(v^{k+2\ell}X_a)$.
For $t<0$ the inverse adjoint has ratio
$\prod_{\ell=t+1}^{0}f(v^{k+2\ell}X_a)$.
Thus $\Ad(U_{a,k,f})^\varepsilon X_m$ has nonnegative integral
Laurent-series coefficients whenever $\varepsilon t\ge0$.
Its constant, or no-jump, coefficient is one.
\end{lemma}

\begin{proof}
Choosing $j$ distinct indices $0\le i_1<\cdots<i_j$ in
\eqref{wall:initial-unit} gives weight
$v^{-j-2\sum i_s}$.  Subtracting $0,1,\ldots,j-1$ from these indices
leaves a weakly increasing nonnegative list.  Its successive
differences have weights $v^{-2},\ldots,v^{-2j}$, which proves
\eqref{wall:unit-coefficient}.  The commutation relation
$X_aX_m=v^{2t}X_mX_a$ proves \eqref{wall:unit-ratio}.  Cancellation of
the two infinite strings leaves the displayed finite strings.
Each factor $1+v^bX_a$, or $(1-v^bX_a)^{-1}$, has nonnegative
coefficients.  Replacing $X_mX_{ja}$ by its centered monomial changes
only the power of $v$.  The case $t=0$ is the identity.
\end{proof}

For later use, the finite formula for the normalized simple unit is
particularly convenient.  If $t=\Omega(p,m)\ne0$ and $a=|t|$, then
\begin{equation}
 \Ad(D_p)^{\operatorname{sign}(t)}X_m
 =\sum_{j=0}^a
 e_j(v^{1-a},v^{3-a},\ldots,v^{a-1})X_{m+jp}.
 \label{wall:binomial}
\end{equation}
Here $e_j$ is the ordinary elementary polynomial in the displayed
$a$ scalar letters.  In particular, when $t=1$ the result is
$X_m+X_{m+p}$, with no additional scalar.

\begin{proposition}[Positive factors on every wall]
\label{wall:factors}
At every generic point of a wall, and to every specified root order,
its element is a product of elementary units on its ray with the
finiteness conditions above.  The same conclusion holds if the
simple incoming data are replaced by allowed products of elementary
units.  If a crossing has the sign of $-V(m)$ on its wall normal,
the resulting monomial jump from $X_m$ has coefficients in
$\N((v^{-1}))$.
\end{proposition}

\begin{proof}
Fix a root-degree bound $d$, and argue inductively in $d$.  For a
given ray $\R_{>0}r$, follow a generic line in $r^\perp$, parallel
to $V(r)$, from far in the positive direction towards the specified
cell.  The initial element is the incoming element just proved.
At the finitely many joints met by the line, all relevant normals
lie in a two-dimensional plane containing $r$.  A normal $s$ with
$\Omega(s,r)=0$ gives a hyperplane parallel to the line, which a
generic offset avoids.  Thus a joint at which a change occurs has
nonzero skew form on its plane.  The offset can also avoid every
intersection involving three independent normals.

Order the positive rays in this plane so that the skew pairing of
an earlier ray with a later ray is positive.  At the incoming end
$V(a)$ of a ray $a$, earlier normals evaluate positively and later
ones negatively.  The local middle element is therefore the
ordered product of its incoming ray elements in this order.  Its
factorization in the reverse order gives its outgoing ray elements.
Theorem~\ref{alg:reordering} applies to precisely this positive
ordering and to the pointed monoid obtained by intersecting the
plane with $C^+$.

There is a small induction issue at degree $d$.  The incoming
element on the ray being followed is already known through degree
$d$ from the previous segment.  On all other incoming rays it is
enough to know positive factorizations through degree $d-1$.
Indeed, a degree-$d$ input term can affect an output term of degree
$d$ only linearly: multiplication with any nonconstant term raises
the degree.  Its linear contribution remains on its own ray.
Unknown degree-$d$ terms on other rays consequently do not affect
the ray under consideration.  Extend their known positive lists
arbitrarily to the required truncation, apply
Theorem~\ref{alg:reordering}, and use uniqueness of ray
factorization.  This proves the assertion at each successive joint
and completes the induction.

For the last assertion, a positive-to-negative crossing has
$\Omega(r,m)>0$ and uses the adjoint of the wall element; a
negative-to-positive crossing has $\Omega(r,m)<0$ and uses its
inverse.  All elementary arguments are positive multiples of $r$,
so Lemma~\ref{wall:unit} applies to every factor.

We spell out why infinite lists still give finite integer
coefficients at a specified energy.  At a fixed root bound there
are finitely many possible arguments, the total multiplicity of
selected nonconstant factors is bounded, and all skew-pairing
shifts between these factors are bounded.  Their parameters have a
common upper bound $K_0$.  If a contribution has energy $E$, no
selected parameter can tend to $-\infty$: the bounded number of
other parameters and the bounded interaction shifts could not
compensate for it.  Every selected parameter is therefore above
an explicit lower threshold depending on the root bound and $E$.
Only finitely many occurrences lie above that threshold.  Their
nonnegative string indices are then bounded as well.  This proves
both finite multiplicity at each energy and a common upper energy
bound.  The corresponding output-list bounds are supplied by
Theorem~\ref{alg:reordering}.  Thus all the asserted coefficients
belong to $\N((v^{-1}))$, rather than merely being formal infinite
sums of nonnegative integers.
\end{proof}

A \emph{section} is a family of values transported by these
consistent wall operations.  Any value with the stipulated support
in one generic chamber determines a section.  Its \emph{incoming
coefficient} is
\begin{equation}
 A_n(F)=[X_n]F(V(n)+\epsilon),
 \label{wall:incoming-coordinate}
\end{equation}
where a sufficiently small generic perturbation is chosen at the
finite order needed for that coefficient.  This does not depend on
the perturbation.  A wall that can separate two such choices has
normal $r$ with $\Omega(r,n)=0$, and its action cannot alter the
$X_n$-coefficient: a possible contributing input $n-jr$ also has
zero pairing with $r$, so that monomial is fixed.

In any fixed chamber $H$, transport adds only positive roots.
Consequently $A_n(F)$ equals $[X_n]F(H)$ plus a linear combination
of coefficients at strictly smaller exponents.  On every finite
interval this is a triangular change of coordinates with diagonal
one.  It defines a unique section $\Theta_m$, supported in
$m+C^+$, such that
\begin{equation}
 A_n(\Theta_m)=\begin{cases}1,&n=m,\\0,&n\ne m.\end{cases}
 \label{wall:theta-definition}
\end{equation}
Its leading monomial is $X_m$ in every chamber.  This incoming-coordinate
description plays the role of the theta functions in the scattering
framework of \cite{GHKK18,DM21}.  The definition and positivity argument
below use only the explicit coefficient recursion in our convention.  More generally,
within any fixed finite union of translates of $C^+$,
\begin{equation}
 F=\sum_n A_n(F)\Theta_n.
 \label{wall:theta-expansion}
\end{equation}
For a requested exponent this is a finite sum, by the interval
property of the free cone.

\begin{lemma}[The positive coefficient recursion]
\label{wall:recursion}
Fix a generic target chamber $H$, represented by $h$, and a target
exponent $n$.  At sufficient finite root order, follow
$h+sV(n)$ from sufficiently large $s$ down to $s=0$.  Let its
successive crossings be $W_i:I_i\longrightarrow I_{i+1}$, with
primitive positive wall normal $p_i$.  Define
\[
 J_i(a,j)=[X_{a+jp_i}]W_i(X_a),\qquad j\ge0.
\]
Then $J_i(a,0)=1$, and for the exponents $a=n-jp_i$ that can
contribute one has $J_i(a,j)\in\N((v^{-1}))$.  Moreover,
\begin{equation}
 [X_n]F(H)=A_n(F)+
 \sum_i\ \sum_{\substack{j>0\\n-jp_i\ \text{in the support bound}}}
 J_i(n-jp_i,j)[X_{n-jp_i}]F(I_i).
 \label{wall:coefficient-recursion}
\end{equation}
All sums are finite at a fixed exponent.  With the simple data
\eqref{wall:initial-unit}, the jump coefficients are effectively
computable at every fixed energy by finite rational arithmetic and
Laurent expansion in $v^{-1}$.  The same holds for the section
coefficients when its initial value is a specified finite Laurent
polynomial, or when its incoming data are finitely specified
rational functions of $v$.
\end{lemma}

\begin{proof}
After rescaling by $s$, the initial point tends to $V(n)$, so its
$X_n$-coefficient is $A_n(F)$.  Hyperplanes whose normals pair
trivially with $n$ are parallel to this line and are avoided by
choosing $h$ generically.  Every actual crossing has the sign of
$-V(n)$; the input $n-jp_i$ has the same pairing with the wall as
$n$.  Proposition~\ref{wall:factors} proves the stated positivity
and the no-jump assertion.  Taking the $X_n$-coefficient in each
crossing identity and summing the successive differences proves
\eqref{wall:coefficient-recursion}.  The exponents in each inner
sum lie in a finite cone interval.  The arrangement contains only
the roots needed for those intervals, so the outer sum is finite
as well.

For effectivity, \eqref{wall:unit-coefficient} gives rational
functions of $v$ for the input coefficients.  The triangular
equations \eqref{wall:incoming-prescription} and
\eqref{wall:three-factors}, followed by inversion and multiplication
at bounded root order, therefore give rational functions of $v$
for every required wall coefficient.  Their Laurent expansions
at $v^{-1}=0$ compute the requested jump coefficients and their
upper energy bounds.  No positive factorization needs to be
found by the algorithm.  A finite polynomial initial value can
be transported by these same finite rational operations; finitely
specified incoming data can instead be used in
\eqref{wall:coefficient-recursion}.  Recursing on its strictly
smaller exponents terminates, since the cone intervals are finite.
At fixed total energy, upper
bounds on the finitely many factor energies also give lower
bounds on each of them.  Hence the required integer coefficient
is obtained by a finite computation.
\end{proof}

\begin{theorem}[Positive sections]
\label{wall:positivity}
Every $\Theta_m$ has coefficients in $\N((v^{-1}))$ in every
chamber.  More generally, the same is true of a section $F$ all
of whose incoming coefficients belong to $\N((v^{-1}))$.
\end{theorem}

\begin{proof}
Fix a target exponent and its finite interval above the support
lower bounds.  Induct on the addition order in this interval,
using \eqref{wall:coefficient-recursion}.  Its initial term is
nonnegative by hypothesis, and every other term is a nonnegative
jump coefficient times a coefficient at a strictly smaller
exponent.  The induction proves positivity.  The energy bounds
in Proposition~\ref{wall:factors} and Lemma~\ref{wall:recursion}
justify each multiplication and sum in the Laurent completion.
Equation~\eqref{wall:theta-definition} gives the special case.
\end{proof}

\begin{corollary}[Positive products]
\label{wall:product}
The structure coefficients in
\[
 \Theta_{m_1}\cdots\Theta_{m_b}
 =\sum_n A_n(\Theta_{m_1}\cdots\Theta_{m_b})\Theta_n
\]
belong to $\N((v^{-1}))$.  Products and expansions are locally
finite at every specified exponent.
\end{corollary}

\begin{proof}
Evaluate the product near $V(n)$ to compute its incoming
$n$-coefficient.  Each factor is positive by
Theorem~\ref{wall:positivity}; the scalar contributed by torus
multiplication is a power of $v$.  For fixed total exponent there
are only finitely many tuples of factor exponents above their
specified lower bounds.  Their Laurent products have finite
coefficients at each energy.  Thus the incoming coefficient is
nonnegative, and \eqref{wall:theta-expansion} proves the assertion.
\end{proof}

\subsection{A finite bound in the mutation direction}

Mutation compatibility is a basic feature of classical cluster theta
functions \cite[Theorem~1.24 and Proposition~3.6]{GHKK18} and their
quantum counterparts \cite[Appendix~A]{DM21}.  We next compare two
different cone completions, keeping the finite-fiber argument needed
for equality of actual coefficients.  From this point
the simple incoming data are exactly
\eqref{wall:initial-unit}.  Fix $p\in C$ and put
\begin{equation}
 \alpha_c=\Omega(p,c),\qquad
 c'=c+[-\alpha_c]_+p\quad(c\ne p),\qquad
 C'=\{-p\}\cup\{c':c\in C\setminus\{p\}\},
 \label{wall:mutated-roots}
\end{equation}
where $[x]_+=\max(x,0)$.  The set $C'$ is again part of a
lattice basis.  Define
\begin{equation}
 S(m)=m-\Omega(p,m)p,
 \qquad S^\vee h=h\circ S^{-1}.
 \label{wall:linear-mutation}
\end{equation}
The map $S$ preserves $\Omega$ and hence acts on the quantum torus
by $X_m\mapsto X_{S(m)}$.  On covectors, the change of chart is
\begin{equation}
 T(h)=\begin{cases}
 h,&h(p)\ge0,\\
 S^\vee h,&h(p)\le0.
 \end{cases}
 \label{wall:piecewise-mutation}
\end{equation}
The two definitions coincide on $p^\perp$.

For a non-pure root, write
\[
 r=kp+u,\qquad u=\sum_{c\ne p}u_cc\ne0,
 \qquad k,u_c\in\N.
\]
Its \emph{non-$p$ degree} is $|u|=\sum_{c\ne p}u_c$.  For
$\sigma\in\{+1,-1\}$ set
\begin{equation}
 M_\sigma(u)=\sum_{c\ne p}[-\sigma\alpha_c]_+u_c.
 \label{wall:fiber-bound-function}
\end{equation}

\begin{lemma}[Finite support in each $p$-fiber]
\label{wall:fiber-bound}
On the open side $\operatorname{sign}h(p)=\sigma$, a non-pure
root $kp+u$ occurring in a wall element or its logarithm satisfies
\begin{equation}
 0\le k\le M_\sigma(u).
 \label{wall:fiber-bound-equation}
\end{equation}
The pure $p$-wall has element $D_p$ everywhere on the cut where
no other wall direction is present.
\end{lemma}

\begin{proof}
Projection to the face $\N p$ is an algebra homomorphism: a sum
of roots with a nonzero non-$p$ component cannot be pure.  At any
point of $p^\perp$, the pure part of the middle factor is thus
the pure part of $g$, which is $D_p$ by
\eqref{wall:incoming-prescription}.  This proves the last
assertion.

We prove \eqref{wall:fiber-bound-equation} by induction on
$|u|$, not on $k+|u|$.  During this induction, an assertion about
an individual coefficient at $kp+u$ is always checked first in a
finite full-degree truncation containing that coefficient.  No
uniform finiteness in $k$ is presumed.

Put $t=\Omega(p,u)=\sum_c\alpha_cu_c$.  If $t\ne0$, call the
side $\sigma=\operatorname{sign}(t)$ upstream.  A line in
$(kp+u)^\perp$, followed from the incoming direction
$V(kp+u)$ to a point on that side, never crosses $p^\perp$:
the evaluation on $p$ along its forward tail is $h(p)+st$ with
$s\ge0$.  If $t=0$, the line is parallel to the cut and the
same argument applies on either side.  At its incoming end the
bound holds, since the only nontrivial incoming rays are simple
ones, with $k=0$.

Consider a joint off the cut along this line.  Every normal
involved has positive non-$p$ degree.  In an ordered
refactorization, a genuinely new contribution of non-$p$ degree
$|u|$ is a product of at least two such terms and therefore
uses only smaller non-$p$ degrees.  A term already of degree
$|u|$ contributes linearly on its own ray, so it is merely
carried to the next segment.  By induction every smaller term
satisfies \eqref{wall:fiber-bound-equation}.  The bound is
additive in $u$: sums of supported exponents are still
supported in
\[
 \{kp+u:0\le k\le M_\sigma(u)\}.
\]
It follows, by the coefficient recursion for ray
factorization, that the bound persists at this joint.  This
proves the upstream assertion for each $k$ separately, through
finitely many joints at the order needed for that $k$.
Only now do we conclude that the entire upstream fiber has
finite support, since the numerical upper bound
$M_\sigma(u)$ is independent of the truncation.

Suppose $t\ne0$ and continue to the other side.  At a generic
transverse joint on $p^\perp\cap u^\perp$, the normals lie
in the plane spanned by $p$ and $u$.  Their non-pure
projections are positive multiples of the same ray through
$u$.  Their pairings with $p$ are consequently nonzero and
have the common sign $\sigma$.  Such genericity is imposed
only through the finite non-$p$ degrees currently needed;
the upstream bounds just proved make the corresponding
full-degree arrangement finite.

Separate the pure factor $D_p$ from the incoming ordered
product of non-pure factors.  If $\sigma=+1$, $p$ is first
in the incoming order, and refactorization conjugates the
non-pure product by $D_p$.  If $\sigma=-1$, $p$ is last,
and it conjugates that product by $D_p^{-1}$.  Formula
\eqref{wall:binomial} shows that an input monomial $X_{kp+w}$
acquires only additional exponents $jp$ with
\[
 0\le j\le |\Omega(p,w)|.
\]
All relevant $w$ have the same pairing sign.  On this ray,
\[
 M_{-\sigma}(w)=M_\sigma(w)+|\Omega(p,w)|.
\]
Thus the whole conjugated non-pure product has the downstream
bound.  So do its ray factors: coefficientwise factorization
cannot introduce an exponent outside an additive support
monoid containing the product's support.  Explicitly, for the
first putative such exponent the defining coefficient is zero
in the product, and every term subtracted from it is a sum of
already permitted exponents, a contradiction.

At the current degree, the coefficients needed at the cut
are either of smaller non-$p$ degree, or already established
upstream coefficients of the current degree.  This avoids a
circular appeal to downstream finiteness.  Further joints
off the cut are treated by the same smaller-degree argument
as before.  Finally, the support monoid is additive, so its
bound is preserved by both logarithm and exponential.  The
statement holds for wall elements and their logarithms.
\end{proof}

The numerical bound also controls the change of completion.
Let $L_+=\id$ and $L_-=S$.  A direct calculation gives
\begin{equation}
 L_\sigma(kp+u)
 =\sum_{c\ne p}u_cc'+(M_\sigma(u)-k)(-p).
 \label{wall:root-image}
\end{equation}
Thus every supported root off the cut maps into $C'^+$.
Put $B=\max_{c\ne p}|\alpha_c|$, with $B=0$ if the index
set is empty.  Since $M_\sigma(u)\le B|u|$, a root of new
degree at most $d$ comes from one with $|u|\le d$ and old
degree at most $(B+1)d$.  In particular, the pushed wall
elements and their cells at bounded new order are determined
by finite old data.

\subsection{Mutation and equality of actual coefficients}

\begin{theorem}[Change of chart]
\label{wall:mutation}
The wall system for $C$, with simple data $D_c$, is carried
to the wall system for $C'$, with simple data $D_{c'}$, as
follows.  Off the cut $p^\perp$, use the covector map $T$
from \eqref{wall:piecewise-mutation} and act on exponents
by $L_\sigma$ on the side of sign $\sigma$.  On the cut,
replace the pure wall by $D_{-p}$.

Suppose $F$ and $F'$ are sections specified by finite Laurent
polynomials in the respective all-negative chambers, and
these initial polynomials satisfy the rational identity
\begin{equation}
 F=\Ad(D_p)F'.
 \label{wall:initial-comparison}
\end{equation}
For $h(p)>0$ their transported values agree coefficientwise
under the identity; for $h(p)<0$ they agree coefficientwise
under $S$.  In particular,
\begin{equation}
 A_n(F)=A_{\mu_p(n)}(F'),\qquad
 \mu_p(n)=n+[-\Omega(p,n)]_+p.
 \label{wall:incoming-mutation}
\end{equation}
The map $\mu_p:M\to M$ is a bijection.
\end{theorem}

\begin{proof}
First we establish the identity of transformations at the
cut, with the normalization of Lemma~\ref{wall:unit}:
\begin{equation}
 S\circ\Ad(D_p)=\Ad(D_{-p})^{-1}.
 \label{wall:cut-identity}
\end{equation}
For $t=\Omega(p,m)\ge0$, the coefficient of $X_{m+jp}$
in $\Ad(D_p)X_m$ is
\[
 b_{t,j}=e_j(v^{1-t},v^{3-t},\ldots,v^{t-1}).
\]
The displayed letters are invariant under inversion and
have product one, so $b_{t,j}=b_{t,t-j}$.  Applying $S$
sends the exponent $m+jp$ to $m-(t-j)p$.
On the other hand, the finite formula for
$\Ad(D_{-p})^{-1}X_m$ has coefficients $b_{t,j}$ at
$m-jp$.  This proves \eqref{wall:cut-identity} on $X_m$.
When $t<0$, apply the proved identity to $X_{-m}$ and take
its inverse.  Both sides are algebra homomorphisms, so the
identity follows on all rational monomial expressions.

\emph{A common ring for the comparison.}
Set $Z=X_p$, choose a lattice complement $M=\Z p\oplus L$
containing $C\setminus\{p\}$, and allow coefficients in
$K(Z)$ on the monomials $X_\ell$, $\ell\in L$.  Complete
in the free positive monoid generated by the images of
$C\setminus\{p\}$ modulo $p$, allowing a finite union of
translates for section values.  Multiplication is specified
by \eqref{wall:torus} and the rule
\begin{equation}
 X_\ell f(Z)=f(v^{-2\Omega(p,\ell)}Z)X_\ell.
 \label{wall:rational-commutation}
\end{equation}
Consequently a coefficient at each specified non-$p$
exponent is a rational function of $Z$.  Products use
finitely many decompositions of that non-$p$ exponent.
This ring will be used only for comparisons; it does not
identify the two Laurent completions in $Z$.

By Lemma~\ref{wall:fiber-bound}, a non-cut wall element is
polynomial in $Z$ at each fixed non-$p$ degree.  Its
constant term in non-$p$ degree is one.  Its inverse,
computed by the geometric series in positive non-$p$
degree, is therefore polynomial in $Z$ at each such degree
as well.  The same assertions hold with Laurent
polynomials after applying either linear piece of the
chart map.  Conjugation by these elements is thus a
well-defined operation on the common ring.  Pure
conjugation by $D_p^{\pm1}$ fixes $Z$ and sends any
$X_\ell$ to a rational function of $Z$ times $X_\ell$,
by \eqref{wall:unit-ratio}.  It, too, acts on this ring.
The map $S$ is the identity modulo $p$, so it preserves
the non-$p$ degree and acts in the same ring.

We justify importing the old consistency into this ring.
The expansion homomorphism
\[
 K(Z)\longrightarrow K((Z))
\]
at $Z=0$ is injective.  On a monomial input, expansion of
every old crossing operation agrees with its operation
in the old positive completion: its increments lie in
$C^+$, and the pure ratios have their usual power-series
expansions at zero.  At bounded non-$p$ degree, all
non-cut wall elements are controlled by finitely many
old roots by Lemma~\ref{wall:fiber-bound}.  The path may
be made generic for these data.  To test a further
coefficient in its expansion in $Z$, perturb through
the additional finite old order needed for that
coefficient; this does not alter the already specified
projected operations.  Old path independence holds at
each such full order.  Hence a rational difference
between two path operations has zero Laurent
expansion in every non-$p$ degree, and injectivity
forces that difference to be zero.  This proves old
consistency in the common ring, on every monomial and
therefore on its completed modules.

\emph{Consistency and incoming data after mutation.}
Push the non-cut elements by $L_\sigma$ and put
$D_{-p}$ on the pure cut.  Equation~\eqref{wall:root-image}
shows that this is a system in the new cone completion,
finite to every new root order.  On any loop, the
linear identifications on successive open sides
cancel; at a cut crossing their discrepancy is
exactly corrected by \eqref{wall:cut-identity}.  Thus
the pushed path operations are consistent in the
common rational ring.  Expand now at $Z^{-1}=0$ to
obtain their operations in the new cone completion.
Replacing $c$ by $c'=c+[-\alpha_c]_+p$ changes the
coefficient at a fixed projected exponent only by a
Laurent monomial in $Z$ and a scalar power of $v$.
Equation~\eqref{wall:root-image} ensures that every
pushed increment is $C'^+$-positive, so this fiberwise
expansion is precisely the one in the new completion.
The finite-order bounds already proved make this
expansion coefficientwise legitimate.  They are
consistent there as well.  Since the form is
nondegenerate, adjoint actions detect the wall group
elements: for a nonzero first root $r$ in the
logarithm of an element, choose $m$ with
$\Omega(r,m)\ne0$; its commutator with $X_m$ is
nonzero.  Consequently consistency of the actions
also gives consistency of the wall products.

It remains to identify this system by its incoming
prescription.  A symplectic linear map $L$ satisfies
$V(Lr)=L^\vee V(r)$.  Thus on each eventual open side,
the incoming direction of a pushed ray is the image
of its old incoming direction.  For a simple $c\ne p$,
$V(c)(p)=\alpha_c$, so that the relevant piece sends
$c$ to exactly $c'=c+[-\alpha_c]_+p$.  Its unit is
therefore $D_{c'}$.  All other non-pure incoming
logarithmic coefficients vanish.  If an incoming
direction is parallel to the cut, its pairing with
$p$ is zero; the two pieces agree on its exponent,
and either sufficiently small generic offset gives
the same conclusion by the trace argument
\eqref{wall:trace}.  A component with no preimage
wall has zero coefficient.  The pure incoming
element is $D_{-p}$ by construction.  Uniqueness of
\eqref{wall:incoming-prescription} identifies the
pushed system with the system for $C'$.

\emph{Actual coefficients of polynomial sections.}
We now use the finite-polynomial hypothesis in
\eqref{wall:initial-comparison}.  Choose a point on
the cut with $h(c)<-1$ for all $c\ne p$, and take
$h(p)$ very small on either side.  Such covectors
exist because $C$ is part of a basis.  With the
constant $B$ defined after
\eqref{wall:root-image}, choosing
$|h(p)|<(B+1)^{-1}$ ensures that every non-pure
supported root off the cut has negative evaluation:
\[
 h(kp+u)\le B|u|\,|h(p)|-|u|<0.
\]
At the cut itself $h(u)<0$.  Hence the short
crossing between these points meets only the pure
wall.  The negative side is in the old all-negative
chamber.  The positive side is in the new
all-negative chamber, since $h(-p)<0$ and
$h(c')<0$.

On the positive side, the old value is
\[
 \Ad(D_p)^{-1}F=F',
\]
which is the actual finite Laurent polynomial
specified in the new chamber.  On the negative
side the pushed old value is $SF$, while
\eqref{wall:cut-identity} and
\eqref{wall:initial-comparison} give
\[
 SF=\Ad(D_{-p})^{-1}F'.
\]
This too is a finite Laurent polynomial, because
$F$ is one and $S$ is linear.  We have therefore
established equality of actual finite values on
both sides near the cut.

To reach any other chamber on a fixed open side,
transport entirely within that half-space.
Only non-cut walls occur.  Their elements and
inverses have finite Laurent support in the
$p$-fiber at every fixed non-$p$ degree, as proved
above.  Starting from a finite Laurent polynomial,
successive transport consequently retains finite
Laurent support in each such fiber: the required
coefficient uses finitely many decompositions of
the non-$p$ degree and finitely many crossings in
its bounded arrangement.  Applying $L_\sigma$
preserves this property.  The compared fiber is
therefore the same finite Laurent polynomial on
the two sides of the identification.  Its
coefficient equality is independent of whether
it is regarded in the old or the new completion.
This establishes the claimed equality of actual
section coefficients throughout both open
sides.  In particular, the argument never
equates the expansions at zero and infinity
of an arbitrary rational function.

Finally, put $t=\Omega(p,n)$.  The incoming
direction $V(n)$ lies on the side of sign $t$.
The relevant linear piece sends $n$ to $n$ if
$t>0$, and to $S(n)=n-tp$ if $t<0$.  If $t=0$,
both pieces fix $n$, and either generic
perturbation computes the same incoming
coefficient.  This proves
\eqref{wall:incoming-mutation}.  Its label map
preserves $\Omega(p,n)$, so its inverse subtracts
$[-\Omega(p,n)]_+p$.  The label map is thus a
bijection, as asserted.
\end{proof}
